% =====================================================================
%  4.2  Atividade 3.2.2.2 -- Realizar trabalhos técnico-sociais
% =====================================================================
\subsection{Atividade 3.2.2.2 -- Realizar trabalhos técnico-sociais nas comunidades e domicílios}
\label{subsec:tts}

\subsubsection{Objetivo e fundamentação}

Obter o diagnóstico social e sanitário das comunidades e domicílios, definir
as famílias a serem atendidas e validar socialmente as tecnologias propostas.
O trabalho técnico-social será executado por empresa contratada, com base em
Termo de Referência elaborado para esse fim, e compreenderá os inventários
sociais, o levantamento das necessidades sanitárias domiciliares e a
validação social das tecnologias \cite{mop2026}.

\begin{fichaatividade}{Ficha-síntese da Atividade 3.2.2.2}{qua:ficha-3222}
  \campo{Código}{3.2.2.2}
  \campo{Tipo de atividade}{Estudo e mobilização}
  \campo{Forma de execução}{Contratação pública (licitação) de empresa especializada}
  \campo{Fonte de recursos}{BIRD}
  \campo{Responsável}{IAT / IDR-Paraná}
  \campo{Dependências institucionais}{Prefeituras}
  \campo{Prazo estimado}{Ano 1 ao ano 4}
  \campo{Produto esperado}{Diagnóstico}
  \campo{Unidade de medida}{Cadastro domiciliar}
  \campo{Evidência de comprovação}{Relatório final e baixa no Termo de Cooperação}
  \campo{Periodicidade}{Por lote}
  \campo{Validação}{IAT / IDR-Paraná}
\end{fichaatividade}

\subsubsection{Procedimentos}

\paragraph{Fase 1 -- Contratação.}
\begin{enumerate}[label=\alph*)]
  \item \textbf{elaboração do Termo de Referência}: documento técnico que
        define escopo, objetivos, requisitos, prazos e produtos esperados;
  \item \textbf{edital de licitação}: elaboração e publicação do edital com
        as regras de contratação;
  \item \textbf{contratação da empresa}.
\end{enumerate}

\paragraph{Fase 2 -- Mobilização e trabalho de campo.}
\begin{enumerate}[label=\alph*), resume]
  \item \textbf{mobilização e sensibilização} dos municípios, comunidades,
        associações e proprietários: atividades para informar a população
        rural sobre o Programa, estimular a participação comunitária e
        conhecer o território onde os sistemas serão implantados;
  \item \textbf{trabalho técnico-social (busca ativa)}: atividades de
        trabalho social com as comunidades, incluindo capacitação e
        organização comunitária;
  \item \textbf{diagnóstico coletivo das formas de abastecimento de água e de
        esgotamento sanitário}: avaliação das condições existentes, da
        demanda de água, da infraestrutura necessária e da viabilidade dos
        sistemas;
  \item \textbf{levantamento das condições domiciliares}.
\end{enumerate}

\paragraph{Fase 3 -- Enquadramento e adesão.}
\begin{enumerate}[label=\alph*), resume]
  \item \textbf{critérios de enquadramento das famílias}: definição, com
        base em aspectos sociais, territoriais e técnicos, das famílias que
        serão atendidas;
  \item \textbf{termo de adesão}: documento assinado pelas famílias
        beneficiadas, que atesta a adesão e o compromisso com o uso e a
        manutenção adequados do sistema.
\end{enumerate}

\paragraph{Fase 4 -- Consolidação.}
\begin{enumerate}[label=\alph*), resume]
  \item \textbf{compilação e apresentação dos resultados}, que subsidiarão a
        definição das comunidades contempladas e os projetos das
        \atividade{3.2.2.3} e \atividade{3.2.2.4}.
\end{enumerate}

\pendencia{No MOP, a contratação da empresa aparece depois da mobilização e
do diagnóstico, dentro de um item intitulado ``Contratação, elaboração,
execução e acompanhamento do Programa'', cuja descrição inclui a contratação
de projetos e obras e, portanto, vai além do trabalho técnico-social. Nesta
estrutura, a contratação foi reposicionada como Fase~1. Confirmar a
sequência e o escopo.}

\pendencia{A atividade não tem quadro de etapas com prazos, produtos
intermediários e evidências, como as \atividade{3.2.2.1} e
\atividade{3.2.2.3} (Quadros 28 e 29 do MOP). Elaborar um quadro no mesmo
formato.}

\subsubsection{Interfaces}

\begin{itemize}
  \item \atividade{3.2.2.1}: o trabalho técnico-social é uma das etapas do
        modelo de gestão (Fase~3).
  \item \atividade{3.2.2.3} e \atividade{3.2.2.4}: o diagnóstico define as
        comunidades contempladas e a condição de disposição final do esgoto
        de cada domicílio.
\end{itemize}
